\documentclass[12pt]{article}
\usepackage[utf8]{inputenc}
\usepackage{amsmath, amssymb, amsthm}
\usepackage[margin=1in]{geometry}

% --------------------------------------------------------------
% Theorem Environments
% --------------------------------------------------------------
\newtheorem{axiom}{Axiom}
\newtheorem{definition}[axiom]{Definition}
\newtheorem{theorem}[axiom]{Theorem}
\newtheorem{lemma}[axiom]{Lemma}
\newtheorem{corollary}[axiom]{Corollary}
\newtheorem{proposition}[axiom]{Proposition}
\newtheorem{remark}[axiom]{Remark}

% --------------------------------------------------------------
% EMK Primitives and Derived Notation
% --------------------------------------------------------------
\newcommand{\Shunya}{\mathfrak{N}}
\newcommand{\cut}{\kappa}
\newcommand{\selfcut}{\chi}
\newcommand{\seam}{\mathcal{S}}
\newcommand{\eye}{\mathsf{Eye}}
\newcommand{\rec}{\mathsf{R}}
\newcommand{\cost}{\rho}
\newcommand{\dist}{d_{\rec}}
\newcommand{\EMK}{\mathrm{EMK}}
\newcommand{\A}{\mathcal{A}}
\title{EMK from Three Axioms}
\author{Derived from $\Shunya$, $\cut$, $\selfcut$}
\date{}

\usepackage{amsmath}
\DeclareMathOperator{\Fix}{Fix}
\begin{document}

\maketitle

\begin{abstract}
We derive the entire EMK framework from exactly three primitive axioms:
the pre-distinction ground $\Shunya$, the first cut $\cut$, and the self-cut $\selfcut$.
No sets, no numbers, no coordinates, no pre-existing space, no metric,
and no primitive recognition are assumed.
Recognition is derived as $\selfcut$ itself.
Equality is derived as collapsed recognition.
Topology $\tau$ emerges as stability under cuts.
The thermodynamic compass $(T,V,S,P)$ and the $\lambda$-hierarchy appear as
alternating and recursive patterns of $\cut$ and $\selfcut$.
\end{abstract}

\section{The Three Axioms}

\begin{axiom}[A0: Pre-Distinction Ground]
There exists $\Shunya$, the pre-distinction ground.
\[
\Shunya \neq \varnothing
\]
$\Shunya$ contains no allocated information, no distinction, no reference, no object.
It is the latent potential before any cut.
\end{axiom}

\begin{axiom}[A1: First Cut]
There exists an act $\cut$ called the \textbf{first cut}.
\[
\cut: \Shunya \longrightarrow X_\cut
\]
$\cut$ allocates information from $\Shunya$ into appearance.
The set $X_\cut = \{p_0\}$ is the first appearance-residue,
containing exactly one primitive mark $p_0$.
This creates the \textbf{object-dimension} $d_{\text{obj}}$.
\end{axiom}

\begin{axiom}[A2: Self-Cut / Recognition]
There exists an act $\selfcut$ called the \textbf{self-cut}.
\[
\selfcut: \text{trace}(\cut) \longrightarrow X_\selfcut
\]
$\selfcut$ observes that a cut has occurred.
It acts on the trace of $\cut$ (the fact that $p_0$ exists).
This creates the \textbf{observer-dimension} $d_{\text{obs}}$.
\[
\boxed{\selfcut \text{ is recognition.}}
\]
No separate recognition primitive is needed.
\end{axiom}

\section{Derived: Appearance Space $X$}

\begin{definition}[Total Appearance Space]
From A1 and A2, define the total appearance space recursively:
\[
X_0 = X_\cut = \{p_0\}
\]
\[
X_{n+1} = X_n \cup \selfcut(X_n) \cup \cut(X_n)
\]
\[
X = \bigcup_{n=0}^{\infty} X_n
\]
\end{definition}

\begin{proposition}[Initial Appearances]
\[
X = \{p_0,\; \selfcut(p_0),\; \cut(p_0),\; \selfcut(\cut(p_0)),\; \cut(\selfcut(p_0)),\; \ldots\}
\]
At minimum, $X$ contains $p_0$ (the first residue) and $\selfcut(p_0)$ (the observer's reflection).
\end{proposition}

\section{Derived: Recognition as $\selfcut$}

\begin{definition}[Recognition]
Recognition is the act $\selfcut$ itself.
For any appearance $x \in X$, $\selfcut(x)$ is the observer's recognition of $x$.
Two appearances are \textbf{recognized as the same} if applying $\selfcut$ to both
produces indistinguishable observer-states.
\end{definition}

\begin{definition}[Recognition Equivalence]
Define the relation $\rec$ on $X$ by:
\[
x \rec y \quad\Longleftrightarrow\quad \selfcut(x) \text{ and } \selfcut(y) \text{ are observer-indistinguishable}
\]
At the primitive level, before equality exists, this is taken as the definition.
\end{definition}

\begin{theorem}[Reflexivity]
For all $x \in X$, $x \rec x$.
\end{theorem}
\begin{proof}
$\selfcut(x)$ is indistinguishable from itself.
\end{proof}

\begin{theorem}[Symmetry]
If $x \rec y$ then $y \rec x$.
\end{theorem}
\begin{proof}
Indistinguishability is symmetric.
\end{proof}

\begin{theorem}[Transitivity]
If $x \rec y$ and $y \rec z$ then $x \rec z$.
\end{theorem}
\begin{proof}
Indistinguishability is transitive.
\end{proof}

Thus $\rec$ is an equivalence relation, even though equality is not yet defined.

\section{Derived: Equality as Collapsed Recognition}

\begin{definition}[Equality]
Equality is the limit case of recognition where all distinguishing cuts have healed.
\[
x = y \quad\Longleftrightarrow\quad \selfcut^n(x) \rec \selfcut^n(y) \text{ for all } n \in \mathbb{N}
\]
Equivalently, $x$ and $y$ belong to the same recognition class and
further cuts produce no new distinction.
\end{definition}

\begin{corollary}
Equality is a special case of recognition, not primitive.
\[
\boxed{\text{Equality} = \text{collapsed recognition}}
\]
\end{corollary}

\section{Derived: Eye Projection}

\begin{definition}[Eye]
The Eye is the map that collapses appearances to their recognition classes:
\[
\eye: X \longrightarrow X/\rec, \qquad \eye(x) = [x]_{\rec}
\]
\end{definition}

\begin{theorem}[Eye Idempotence]
\[
\eye^2 = \eye
\]
\end{theorem}
\begin{proof}
Applying $\eye$ to a recognition class returns the same class.
\end{proof}

\section{Derived: Seam}

\begin{definition}[Seam of a Cut]
For any cut $C$ (including $\cut$ and composites), the seam is:
\[
\seam_C = \{ x \in X \mid \selfcut(C(x)) \rec \selfcut(x) \}
\]
The seam is where cutting produces no new recognizable distinction.
\end{definition}

\begin{theorem}[Seam as Fixed Point of $\selfcut \circ C$]
\[
\seam_C = \Fix(\selfcut \circ C) \quad \text{up to recognition}
\]
\end{theorem}

\section{Derived: Cut Family $\mathcal{C}$}

\begin{definition}[Cut Family]
The family $\mathcal{C}$ of admissible cuts is generated by:
\begin{enumerate}
    \item The first cut $\cut$
    \item The self-cut $\selfcut$
    \item Composition: if $C_1, C_2 \in \mathcal{C}$, then $C_1 \circ C_2 \in \mathcal{C}$
    \item Recursive application: $\cut$ applied to any appearance in $X$
\end{enumerate}
\[
\mathcal{C} = \langle \cut, \selfcut \rangle
\]
\end{definition}

\section{Derived: Recognition Distance}

\begin{definition}[Cost of a Cut]
For a cut $C \in \mathcal{C}$, define the cost at $x$ as:
\[
\cost_C(x) = \min\{ n \in \mathbb{N} \cup \{\infty\} \mid \selfcut^n(C(x)) \rec \selfcut^n(x) \}
\]
The cost is the minimal number of self-cut applications needed to heal the distinction.
\end{definition}

\begin{definition}[Recognition Distance]
For recognition classes $a, b \in X/\rec$, define:
\[
\dist(a,b) = \inf\{ \cost_C(x) \mid C \in \mathcal{C},\; x \in a,\; C(x) \in b \}
\]
This satisfies:
\[
\dist(a,b) = 0 \iff a = b
\]
\end{definition}

\section{Derived: EMK Topology}

\begin{definition}[EMK-Open Set]
A subset $U \subseteq X$ is EMK-open if it is stable under all admissible cuts:
\[
U \in \tau \quad\Longleftrightarrow\quad \eye(C(U)) \subseteq \eye(U) \quad \forall C \in \mathcal{C}
\]
\end{definition}

\begin{theorem}[Topology Theorem]
$\tau$ forms a topology on $X$ provided $\eye$ preserves unions and finite intersections.
\end{theorem}
\begin{proof}[Sketch]
$\varnothing, X \in \tau$ by vacuity.
Union and finite intersection closure follow from the preservation properties of $\eye$.
\end{proof}

\section{Derived: Flow and Time from $\selfcut$}

\begin{definition}[Flow Generator]
The self-cut $\selfcut$ generates a flow by repeated application:
\[
\Phi_t(x) = \selfcut^{t}(x)
\]
where $t$ is a counting index (not a real number yet).
\end{definition}

\begin{definition}[Time Direction]
The observer, created by $\selfcut$, experiences the sequence:
\[
x,\; \selfcut(x),\; \selfcut^2(x),\; \ldots
\]
This ordered sequence is \textbf{time}. No external clock is assumed.
\end{definition}

\begin{theorem}[Time Emerges from $\selfcut$]
\[
\boxed{\text{Time} = \text{iterated self-cut}}
\]
\end{theorem}

\section{Derived: Helical Torus as Dynamic Orbit}

\begin{definition}[Helical Torus Orbit]
For a seed appearance $x_0$, define the dynamic helical torus:
\[
\mathcal{H}(x_0) = \{ \selfcut^n(\cut^m(x_0)) \mid n,m \in \mathbb{N} \}
\]
Alternating $\cut$ and $\selfcut$ produces a spiral structure.
Without $\selfcut$, the orbit is static (frozen helix).
With $\selfcut$, the orbit becomes a flow.
\end{definition}

\section{Derived: TVSP Compass from Alternating Cuts}

\begin{definition}[Four Directions]
Alternate $\cut$ and $\selfcut$ in a 4-cycle:
\[
\cut \longrightarrow \selfcut \longrightarrow \cut \longrightarrow \selfcut \longrightarrow \cut
\]
Interpret these as the four thermodynamic directions:
\[
\boxed{T \;\xrightarrow{\cut}\; V \;\xrightarrow{\selfcut}\; S \;\xrightarrow{\cut}\; P \;\xrightarrow{\selfcut}\; T}
\]
\end{definition}

\begin{theorem}[TVSP Compass Emergence]
The cyclic order $T \to V \to S \to P \to T$ is the minimal closed pattern
generated by alternating $\cut$ and $\selfcut$.
\end{theorem}

\section{Derived: $\lambda$-Layers as Recursive Cuts}

\begin{definition}[Recursive Response]
Apply the same cut pattern to the response variables:
\[
\lambda^{(1)}_p = \selfcut(\cut(T,V,S,P))
\]
\[
\lambda^{(2)} = \selfcut(\cut(\lambda^{(1)}))
\]
In general:
\[
\lambda^{(n+1)} = \selfcut(\cut(\lambda^{(n)}))
\]
\end{definition}

\begin{theorem}[Hierarchy of Layers]
Each layer $\mathcal{L}_n$ has a different seam structure because
what constitutes a "recognizable distinction" changes at each recursive level.
\[
\mathcal{L}_0 \neq_{\tau} \mathcal{L}_1 \neq_{\tau} \mathcal{L}_2 \neq_{\tau} \cdots
\]
\end{theorem}

\section{Derived: Real Numbers from Infinite Recursion}

\begin{definition}[Cut-Counting]
Let $t \in \mathbb{R}$ be represented by an infinite sequence of alternating $\cut$ and $\selfcut$:
\[
t \sim (\cut, \selfcut, \cut, \selfcut, \ldots)
\]
Rational numbers appear as periodic sequences.
Irrational numbers appear as non-periodic sequences.
\end{definition}

\begin{theorem}[Real Numbers are Cut-Orbits]
\[
\boxed{\mathbb{R} \cong \{ \text{infinite alternating cut-selfcut sequences} \} / \sim_{\rec}}
\]
The real line is the recognition-collapse of the infinite cut tower.
\end{theorem}

\section{Conclusion: Three Axioms Suffice}

From exactly three axioms $(\Shunya, \cut, \selfcut)$ we have derived:

\begin{enumerate}
    \item Appearance space $X$
    \item Recognition $\rec = \selfcut$
    \item Equality as collapsed recognition
    \item Eye projection $\eye$
    \item Seam $\seam_C$
    \item Cut family $\mathcal{C}$
    \item Recognition distance $\dist$
    \item EMK topology $\tau$
    \item Flow and time from iterated $\selfcut$
    \item Helical torus as dynamic orbit
    \item TVSP compass $(T,V,S,P)$ from alternating cuts
    \item $\lambda$-hierarchy from recursive cuts
    \item Real numbers from infinite cut sequences
\end{enumerate}

\[
\boxed{
\text{No hidden assumptions. No primitive recognition. No primitive equality.}
}
\]

\[
\boxed{
\text{Three axioms. Everything derived.}
}
\]
\section{Three Primitives: The Only Source}

\begin{axiom}[A0: Pre-Distinction Ground]
There exists \(\mathfrak N\), the pre-distinction ground (Śūnya).
\[
\mathfrak N \neq \varnothing
\]
\(\mathfrak N\) contains no allocated information, no distinction, no reference, no object.
It is the latent potential before any cut.
\end{axiom}

\begin{axiom}[A1: First Cut / Observation]
There exists an act \(\kappa\) called the \textbf{first cut}.
\[
\kappa: \mathfrak N \longrightarrow X_\kappa
\]
\(\kappa\) allocates information from \(\mathfrak N\) into appearance.
The set \(X_\kappa = \{p_0\}\) is the first appearance-residue,
containing exactly one primitive mark \(p_0\).
This creates the \textbf{object-dimension} \(d_{\text{obj}}\).
\end{axiom}

\begin{axiom}[A2: Self-Cut / Recognition]
There exists an act \(\chi\) called the \textbf{self-cut}.
\[
\chi: \text{trace}(\kappa) \longrightarrow X_\chi
\]
\(\chi\) observes that a cut has occurred. It acts on the trace of \(\kappa\) (the fact that \(p_0\) exists).
This creates the \textbf{observer-dimension} \(d_{\text{obs}}\).
\[
\boxed{\chi \text{ is recognition.}}
\]
No separate recognition primitive is needed.
\end{axiom}

\section{Derived: Appearance Space \(X\)}

\begin{definition}[Total Appearance Space]
From A1 and A2, define the total appearance space recursively:
\[
X_0 = X_\kappa = \{p_0\}
\]
\[
X_{n+1} = X_n \cup \kappa(X_n) \cup \chi(X_n)
\]
\[
X = \bigcup_{n=0}^{\infty} X_n
\]
Every appearance is a finite sequence of \(\kappa\) and \(\chi\) applied to the seed.
\end{definition}

\section{Derived: Recognition}

\begin{definition}[Recognition]
Recognition is the act \(\chi\) itself.
For any appearance \(x \in X\), \(\chi(x)\) is the observer's recognition of \(x\).
Two appearances are recognized as the same if applying \(\chi\) to both produces indistinguishable observer-states.
\[
\boxed{\text{Recognition} = \chi}
\]
\end{definition}

\section{Derived: Equality as Collapsed Recognition}

\begin{definition}[Equality]
Equality is the limit case of recognition where all distinguishing cuts have healed.
\[
x = y \quad\Longleftrightarrow\quad \chi^n(x) \sim \chi^n(y) \text{ for all } n \in \mathbb{N}
\]
Equivalently, \(x\) and \(y\) belong to the same recognition class and further cuts produce no new distinction.
\[
\boxed{\text{Equality} = \text{collapsed recognition}}
\]
\end{definition}

\section{Derived: Recognition Relation and Eye}

\begin{definition}[Recognition Relation]
Define \(R \subseteq X \times X\) by:
\[
x \sim_R y \quad\Longleftrightarrow\quad \chi(x) \text{ and } \chi(y) \text{ are observer-indistinguishable}
\]
\(R\) is an equivalence relation (reflexive, symmetric, transitive).
\end{definition}

\begin{definition}[Recognition Class]
For \(x \in X\), define its recognition class by:
\[
[x]_R = \{ y \in X : y \sim_R x \}
\]
\end{definition}

\begin{definition}[Eye Map]
The Eye is the recognition collapse projection:
\[
E: X \to X/R, \qquad E(x) = [x]_R
\]
\end{definition}

\begin{theorem}[Eye Idempotence]
\[
E^2 = E
\]
\end{theorem}

\section{Derived: Cut Family \(\mathcal{C}\)}

\begin{definition}[Cut Family]
The family \(\mathcal{C}\) of admissible cuts is generated by:
\begin{enumerate}
    \item The first cut \(\kappa\)
    \item The self-cut \(\chi\)
    \item Composition: if \(C_1, C_2 \in \mathcal{C}\), then \(C_1 \circ C_2 \in \mathcal{C}\)
    \item Recursive application: \(\kappa\) applied to any appearance in \(X\)
\end{enumerate}
\[
\mathcal{C} = \langle \kappa, \chi \rangle
\]
Thus \(\mathcal{C}\) is derived, not primitive.
\end{definition}

\section{Derived: Recognition Distance}

\begin{definition}[Cost of a Cut]
For a cut \(C \in \mathcal{C}\) and \(x \in X\), define the cost:
\[
\rho_C(x) = \min\{ n \in \mathbb{N} \cup \{\infty\} \mid \chi^n(C(x)) \sim_R \chi^n(x) \}
\]
The cost is the minimal number of self-cuts needed to heal the distinction.
\end{definition}

\begin{definition}[Recognition Distance]
For recognition classes \(a, b \in X/R\), define:
\[
d_R(a,b) = \inf\{ \rho_C(x) \mid C \in \mathcal{C},\; x \in a,\; C(x) \in b \}
\]
This satisfies \(d_R(a,b) = 0 \iff a = b\).
\end{definition}

\section{Derived: Seam}

\begin{definition}[Seam of a Cut]
For a cut \(C \in \mathcal{C}\), define its seam by:
\[
S_C = \{ x \in X : \chi(C(x)) \sim_R \chi(x) \}
\]
The seam is the locus where the cut creates no recognizable distinction.
\end{definition}

\begin{theorem}[Seam Zero-Cost Theorem]
\[
x \in S_C \quad\Longleftrightarrow\quad \rho_C(x) = 0
\]
Hence \(S_C = \rho_C^{-1}(0)\).
\end{theorem}

\section{Derived: EMK Topology}

\begin{definition}[EMK-Open Set]
A subset \(U \subseteq X\) is EMK-open if it is stable under all cuts:
\[
U \in \tau_E \quad\Longleftrightarrow\quad E(C(U)) \subseteq E(U) \quad \forall C \in \mathcal{C}
\]
\end{definition}

\begin{theorem}[EMK Topology Theorem]
\(\tau_E\) forms a topology on \(X\) provided \(E\) preserves unions and finite intersections.
\end{theorem}
% ============================================================
% DERIVATION OF PATH EQUALITY FROM κ AND χ
% ============================================================

\section{Angles from Alternating Cuts}

\begin{definition}[4-Cycle Property]
From Axioms A1 and A2, consider the alternating sequence:
\[
\kappa \;\longrightarrow\; \chi \;\longrightarrow\; \kappa \;\longrightarrow\; \chi \;\longrightarrow\; \kappa
\]
After four steps, the composition returns to the starting recognition class up to equivalence:
\[
\chi \circ \kappa \circ \chi \circ \kappa \;\sim_{\rec}\; \text{id}
\]
This is the \textbf{4-cycle property} of alternating cut and self-cut.
\end{definition}

\begin{definition}[Angle Parameterization]
For each recognition dimension \(j\), let \(\theta_j \in [0, 2\pi)\) parameterize the \(j\)-th 4-cycle.
Moving along \(\theta_j\) applies the alternating sequence continuously.
The seam phase \(\gamma\) parameterizes the diagonal where all \(\theta_j\) align.
\end{definition}

\begin{lemma}[Cycle Orthogonality]
Distinct recognition dimensions \(j \neq k\) correspond to independent 4-cycles.
Their generators commute up to recognition:
\[
[\theta_j, \theta_k] = 0 \quad \text{in the recognition quotient}
\]
\end{lemma}

\section{Derived: Path Equality Along Angles}

\begin{definition}[Recognition-Invariant Function]
A function \(f: X \to \mathbb{C}\) is \textbf{recognition-invariant} along the \(\theta_j\)-cycle if:
\[
f(\theta_j + d\theta_j) = f(\theta_j) \cdot e^{i\,d\theta_j}
\]
up to recognition equivalence. This means the phase accumulates uniformly.
\end{definition}

\begin{theorem}[Path Equality in One Dimension]
For any recognition-invariant function \(f\) and any two same-endpoint paths \(\eta_1, \eta_2\) in the \(\theta_j\)-cycle:
\[
\int_{\eta_1} \frac{df}{f} = \int_{\eta_2} \frac{df}{f}
\]
\end{theorem}

\begin{proof}
From the recognition-invariance condition:
\[
\frac{df}{f} = i\,d\theta_j
\]
The right-hand side depends only on the endpoints (the total change in \(\theta_j\)), not on the path.
Therefore the integral is path-independent.
\end{proof}

\begin{theorem}[Path Equality in All Dimensions]
For a function \(f\) that is recognition-invariant along every \(\theta_j\)-cycle \((j = 1, 2, 3, \ldots)\):
\[
\frac{df}{f} = i\sum_{j=1}^{\infty} d\theta_j + i\,d\gamma
\]
where \(\gamma\) is the seam phase. Consequently, for any two same-endpoint paths in the infinite torus:
\[
\int_{\eta_1^{(j)}} \frac{df}{f} = \int_{\eta_2^{(j)}} \frac{df}{f} \qquad \forall j
\]
\end{theorem}

\begin{proof}
By the orthogonality lemma, each \(\theta_j\) contributes independently.
Summing over all dimensions gives the total differential.
Path independence in each direction follows from the one-dimensional theorem.
\end{proof}

\section{Derived: Master Seam Coordinate}

\begin{definition}[Master Seam Coordinate]
Define the \textbf{master seam coordinate}:
\[
\Xi_{\mathrm{EMK}} = \gamma + \sum_{j=1}^{\infty} \theta_j
\]
where the sum is formal — only finitely many terms are nonzero at any finite recognition depth.
The infinite sum is defined as the limit toward the Eye.
\end{definition}

\begin{theorem}[EMK Analytic Form]
If \(f\) satisfies path equality in all angular directions and the seam phase \(\gamma\) follows the same orientation, then:
\[
f = C \exp\left(i \Xi_{\EMK}\right) = C \exp\left(i\left(\gamma + \sum_{j=1}^{\infty} \theta_j\right)\right)
\]
More generally, any \textbf{weak cut-analytic function} has the form:
\[
f(\Theta, \gamma) = F(e^{i\Xi_{\EMK}})
\]
where \(F\) is an arbitrary function of the master seam phase.
\end{theorem}

\begin{proof}
Integrate \(\frac{df}{f} = i\,d\Xi_{\EMK}\) to obtain \(\ln f = i\Xi_{\EMK} + \text{constant}\).
Exponentiating gives the result.
The weak form follows because any function of \(e^{i\Xi_{\EMK}}\) also satisfies path equality.
\end{proof}

\section{Derived: Curvature as Failure of Path Equality}

\begin{definition}[Cut-Order Distinguishability]
For two cuts \(C_i, C_j \in \mathcal{C}\), define:
\[
\Delta_{ij}(x) = d_R\big(E(C_i C_j(x)),\; E(C_j C_i(x))\big)
\]
where \(d_R\) is the recognition distance.
\end{definition}

\begin{theorem}[Curvature Obstructs Path Equality]
If \(\Delta_{ij}(x) \neq 0\), then the order of cuts matters.
Consequently, the angular paths in the corresponding dimensions are not independent:
\[
\int \frac{df}{f} \text{ depends on path order when } \Delta_{ij} \neq 0
\]
Thus curvature is exactly the obstruction to simultaneous path equality in all dimensions.
\end{theorem}

\begin{corollary}[Flatness = Full Path Equality]
When \(\Delta_{ij} = 0\) for all \(i,j\), then:
\[
\frac{df}{f} = i\sum_j d\theta_j + i d\gamma
\]
holds globally, and all angular integrals are path-independent.
This is the \textbf{flat EMK geometry} (equilibrium limit).
\end{corollary}

\begin{remark}[Onsager as Flat Limit]
The standard Onsager reciprocity relations hold exactly when \(\Delta_{ij} = 0\).
Thus Onsager is the flat limit of EMK geometry.
Non-equilibrium corresponds to non-zero curvature (\(\Delta_{ij} \neq 0\)).
\end{remark}

\section{Derived: EMK Analytic Class}

\begin{definition}[EMK Analytic Class]
The \textbf{EMK analytic class} \(\mathcal{A}_{\EMK}\) is the set of functions \(f: X \to \mathbb{C}\) satisfying:
\[
\int_{\eta_1^{(j)}} \frac{df}{f} = \int_{\eta_2^{(j)}} \frac{df}{f} \qquad \forall j,\; \forall \text{ same-endpoint paths } \eta_1^{(j)}, \eta_2^{(j)}
\]
Equivalently:
\[
\mathbb{A}_{\EMK} = \left\{ f : f = F(e^{i\Xi_{\EMK}}) \text{ for some } F \right\}
\]
\end{definition}

\begin{theorem}[Global Angular Analyticity is Derived]
The condition of global angular analyticity — path independence in all angular directions — is \textbf{derived} from the three axioms \((\mathfrak N, \kappa, \chi)\) via:
\begin{enumerate}
    \item The 4-cycle property of alternating \(\kappa\) and \(\chi\)
    \item The orthogonality of distinct recognition dimensions
    \item Recognition-invariance of analytic functions along cycles
\end{enumerate}
Thus it is a theorem, not an additional axiom.
\end{theorem}
\section{Derivation of EMK Analyticity from Three Axioms}

\begin{theorem}[4-Cycle Property]
From \(\kappa\) and \(\chi\), the composition \(\chi \circ \kappa \circ \chi \circ \kappa\) returns to the starting recognition class.
Thus \((\chi \circ \kappa)^2 \sim_{\rec} \text{id}\).
\end{theorem}

\begin{proof}
\(\kappa\) allocates distinction. \(\chi\) observes it. After two allocations and two observations, the observer has completed a full cycle of distinction and reflection. By definition of recognition as \(\chi\), the start and end states are recognized as the same.
\end{proof}

\begin{definition}[Angles]
Each independent 4-cycle generates a parameter \(\theta_j \in [0, 2\pi)\). The collection \(\Theta = (\theta_1, \theta_2, \ldots)\) lives in \(\mathbb{T}^\infty\).
\end{definition}

\begin{definition}[Recognition-Invariant Function]
A function \(f\) is recognition-invariant along \(\theta_j\) if:
\[
f(\theta_j + d\theta_j) = f(\theta_j) e^{i d\theta_j}
\]
This follows from \(\chi\) being the recognition act and the cycle being closed.
\end{definition}

\begin{theorem}[Path Equality]
For any recognition-invariant \(f\) and any same-endpoint paths \(\eta_1, \eta_2\) in \(\theta_j\):
\[
\int_{\eta_1} \frac{df}{f} = \int_{\eta_2} \frac{df}{f}
\]
\end{theorem}

\begin{proof}
From invariance, \(\frac{df}{f} = i d\theta_j\). The integral depends only on \(\Delta \theta_j\), not on the path.
\end{proof}

\begin{theorem}[Master Seam Coordinate]
Summing over all dimensions:
\[
\frac{df}{f} = i\left(\sum_j d\theta_j + d\gamma\right)
\]
where \(\gamma\) is the seam phase. Define \(\Xi = \gamma + \sum_j \theta_j\). Then:
\[
f = C e^{i\Xi}
\]
\end{theorem}

\begin{corollary}[EMK Analytic Class]
\[
\mathcal{A}_{\EMK} = \{ f : f = F(e^{i\Xi}) \}
\]
\end{corollary}

\begin{remark}
Global angular analyticity — path independence in all \(\theta_j\) simultaneously — is therefore \textbf{derived} from \((\mathfrak N, \kappa, \chi)\), not an additional axiom.
\end{remark}
\section{Theorem: \(\tau_E\) is the Initial Topology}

\begin{theorem}[Characterization of EMK Topology]
Let \(\mathcal{S} = \{ E^{-1}(V) : V \subseteq X/R \text{ is cut-invariant} \} \cup \{ f^{-1}(W) : f \in \A_{\EMK},\ W \subseteq \mathbb{C} \text{ open} \}\).
Then the topology generated by \(\mathcal{S}\) as a subbasis equals \(\tau_E\).
\end{theorem}

\begin{proof}
(1) \(\mathcal{S} \subseteq \tau_E\):
\begin{itemize}
    \item For \(f \in \A_{\EMK}\), \(f = \tilde{f} \circ E\) with \(\tilde{f}\) recognition-invariant.
    Then \(f^{-1}(W) = E^{-1}(\tilde{f}^{-1}(W))\).
    For any cut \(C\), \(\tilde{f}^{-1}(W)\) is cut-invariant because \(\tilde{f}(C(a)) = \tilde{f}(a)\).
    Hence \(E(C(f^{-1}(W))) \subseteq \tilde{f}^{-1}(W) = E(f^{-1}(W))\).
    Thus \(f^{-1}(W) \in \tau_E\).
    \item For any cut-invariant \(V \subseteq X/R\), \(E^{-1}(V) \in \tau_E\) by definition of cut-invariance.
\end{itemize}
Thus \(\mathcal{S} \subseteq \tau_E\), and since \(\tau_E\) is a topology, \(\langle \mathcal{S} \rangle \subseteq \tau_E\).

(2) \(\tau_E \subseteq \langle \mathcal{S} \rangle\):
Take any \(U \in \tau_E\). For each \(x \in U\), let \(a = E(x)\).
Because \(U\) is cut-stable, the entire fiber \(E^{-1}(a)\) is contained in \(U\).
Moreover, for any cut \(C\), \(C(a) \in E(U)\).
The set \(E(U)\) is cut-invariant in \(X/R\).
By the derivation of path equality, cut-invariant subsets of \(X/R\) are precisely those that are unions of level sets of \(\tilde{f}\) for \(f \in \A_{\EMK}\).
Thus \(E(U) = \bigcup_{\alpha} \tilde{f}_\alpha^{-1}(W_\alpha)\) for some family.
Then \(U = \bigcup_{\alpha} E^{-1}(\tilde{f}_\alpha^{-1}(W_\alpha)) = \bigcup_{\alpha} f_\alpha^{-1}(W_\alpha)\).
Each \(f_\alpha^{-1}(W_\alpha)\) is a finite intersection of subbasis elements (actually a single subbasis element).
Hence \(U \in \langle \mathcal{S} \rangle\).

Therefore \(\tau_E = \langle \mathcal{S} \rangle\).
\end{proof}

\begin{corollary}
\(\tau_E\) is the initial topology generated by the family \(\{E\} \cup \A_{\EMK}\).
That is, it is the coarsest topology making \(E\) continuous (with respect to the cut-invariant topology on \(X/R\)) and all \(f \in \A_{\EMK}\) continuous.
\end{corollary}
% ============================================================
% TOPOLOGY OF TVSP SPACE FROM EMK TOPOLOGY τ_E
% ============================================================

\section{TVSP Space as a Recognition Quotient}

\begin{definition}[TVSP Equivalence]
From the 4-cycle property derived in Section [X], the alternating sequence
\[
\kappa \circ \chi \circ \kappa \circ \chi \sim_{\rec} \text{id}
\]
generates an equivalence relation \(\sim_{\text{cycle}}\) on \(X/R\):
\[
T \sim V \sim S \sim P \sim T
\]
where \(T, V, S, P\) are the recognition classes corresponding to the four directions.
\end{definition}

\begin{definition}[TVSP Space]
Define the TVSP space as the quotient:
\[
X_{\text{TVSP}} := (X/R) / \sim_{\text{cycle}}
\]
Equivalently, \(X_{\text{TVSP}}\) is the set of equivalence classes of the 4-cycle action.
It has exactly one point when considered as a set of cycles, but better: it is the 4-element set \(\{T, V, S, P\}\) with the cyclic identification. For topology, we treat it as the set \(\{T, V, S, P\}\) before identification, then take quotient topology.
\end{definition}

\begin{remark}
The TVSP space is not a manifold in the classical sense. It is a \textbf{compass} — a directed cyclic structure. Its topological properties are derived from the cut-stability condition.
\end{remark}

\section{Induced Topology on TVSP Space}

Let \(q: X \to X_{\text{TVSP}}\) be the quotient map that first collapses by recognition \(E: X \to X/R\), then by the cycle equivalence.

\begin{definition}[Quotient Topology on TVSP]
A subset \(U \subseteq X_{\text{TVSP}}\) is open if and only if its preimage \(q^{-1}(U) \subseteq X\) is open in \(\tau_E\).
\end{definition}

\begin{lemma}[Factorization]
The quotient map factors as \(q = \bar{q} \circ E\), where:
\[
X \xrightarrow{E} X/R \xrightarrow{\bar{q}} X_{\text{TVSP}}
\]
Thus \(q^{-1}(U) = E^{-1}(\bar{q}^{-1}(U))\).
\end{lemma}

\begin{lemma}[Cut-Invariance Condition]
From the definition of \(\tau_E\), a set \(U \subseteq X_{\text{TVSP}}\) is open iff \(\bar{q}^{-1}(U) \subseteq X/R\) is cut-invariant:
\[
\forall C \in \mathcal{C}, \quad C(\bar{q}^{-1}(U)) \subseteq \bar{q}^{-1}(U)
\]
where the action of cuts on \(X/R\) is induced by \(C([x]) = [C(x)]\).
\end{lemma}

\section{Action of Cuts on TVSP Classes}

From the 4-cycle \(T \xrightarrow{\kappa} V \xrightarrow{\chi} S \xrightarrow{\kappa} P \xrightarrow{\chi} T\), the cuts act as:

\[
\begin{array}{c|cccc}
C & T & V & S & P \\
\hline
\kappa & V & - & P & - \\
\chi & - & S & - & T \\
\end{array}
\]

Where \(-\) indicates that the cut is not directly defined on that class, but compositions extend the action.

\begin{lemma}[Full Cyclic Action]
The group generated by \(\kappa\) and \(\chi\) acts transitively on \(\{T, V, S, P\}\) as the cyclic group \(\mathbb{Z}_4\). Explicitly:
\[
\kappa(T) = V,\quad \chi(V) = S,\quad \kappa(S) = P,\quad \chi(P) = T
\]
and all other actions are obtained by composition.
\end{lemma}

\begin{proof}
From the definitions:
\[
\kappa(T) = V \quad \text{(by the cycle)}
\]
\[
\chi(V) = S \quad \text{(by the cycle)}
\]
\[
\kappa(S) = P \quad \text{(by the cycle)}
\]
\[
\chi(P) = T \quad \text{(by the cycle)}
\]
Compositions give the full cyclic permutation of order 4.
\end{proof}

\section{Classification of Cut-Invariant Subsets}

\begin{theorem}[Cut-Invariant Subsets of \(\{T, V, S, P\}\)]
A subset \(A \subseteq \{T, V, S, P\}\) is invariant under the action of all cuts \(C \in \mathcal{C}\) if and only if \(A = \varnothing\) or \(A = \{T, V, S, P\}\).
\end{theorem}

\begin{proof}
Let \(A\) be non-empty. Pick any element, say \(T \in A\).
- By \(\kappa(T) = V\), we have \(V \in A\).
- By \(\chi(V) = S\), we have \(S \in A\).
- By \(\kappa(S) = P\), we have \(P \in A\).
Thus \(A\) contains all four elements.
The same argument starting from any element cycles through all four.
Therefore the only non-empty invariant subset is the whole set.
The empty set is trivially invariant.
\end{proof}

\section{Topology of TVSP Space}

\begin{theorem}[Indiscrete Topology]
The quotient topology on \(X_{\text{TVSP}}\) induced by \(\tau_E\) is the indiscrete topology:
\[
\tau_{\text{TVSP}} = \{ \varnothing,\ X_{\text{TVSP}} \}
\]
\end{theorem}

\begin{proof}
Take any non-empty \(U \subseteq X_{\text{TVSP}}\). Its preimage in \(X/R\) is \(\bar{q}^{-1}(U)\). Since \(\bar{q}\) identifies all four points under the cycle, \(\bar{q}^{-1}(U)\) is either \(\varnothing\) or the entire set \(\{T, V, S, P\}\) depending on whether \(U\) is empty or not.
If \(U \neq \varnothing\), then \(\bar{q}^{-1}(U) = \{T, V, S, P\}\).
By the theorem above, \(\{T, V, S, P\}\) is cut-invariant.
Hence \(U\) is open.
Thus every non-empty subset of \(X_{\text{TVSP}}\) is open, and the empty set is open by definition.
Therefore \(\tau_{\text{TVSP}}\) is the indiscrete topology.
\end{proof}

\section{Interpretation}

\begin{remark}[Topological Triviality of TVSP]
The TVSP space, as a topological space, carries no non-trivial information. Its only open sets are the empty set and the whole space.
This does \textbf{not} mean the TVSP compass is trivial. The information resides in:
\begin{enumerate}
    \item The \textbf{cut actions} \(\kappa\) and \(\chi\) (the arrows)
    \item The \textbf{recognition classes} \(T, V, S, P\) as distinct points before identification
    \item The \textbf{cyclic order} \(T \to V \to S \to P \to T\)
\end{enumerate}
The topology is trivial because the 4-cycle forces all points to be indistinguishable under cuts. This is a feature, not a bug: the compass is a \textbf{directed graph}, not a space with non-trivial open sets.
\end{remark}

\begin{corollary}[Comparison with Classical Thermodynamics]
In classical thermodynamics, the state space \((T, V, S, P)\) is typically given the standard manifold topology, where open sets are coordinate balls.
In EMK, the topology induced by \(\tau_E\) is indiscrete.
This is not a contradiction because:
\begin{enumerate}
    \item The classical topology is a \textbf{representation choice}, not derived from first principles.
    \item EMK topology is derived from cut-stability, which encodes distinguishability under observation.
    \item The indiscrete topology reflects the fact that under the 4-cycle, all thermodynamic states are connected by cuts — none is more "open" than another.
\end{enumerate}
\end{corollary}

\section{Summary}

\[
\boxed{
\begin{array}{c}
X_{\text{TVSP}} \text{ has the indiscrete topology } \{\varnothing, X_{\text{TVSP}}\}. \\
\text{The thermodynamic compass is a directed cyclic graph, not a manifold.}
\end{array}
}
\]
% ============================================================
% DIRECTED GRAPH STRUCTURE OF TVSP COMPASS
% ============================================================

\section{Categorical Structure vs. Topological Structure}

The previous section showed that the TVSP space \(X_{\text{TVSP}}\) inherits the indiscrete topology from \(\tau_E\). This means that open sets carry no information. However, the TVSP compass is not trivial — its structure is encoded in the \textbf{directed arrows} (cuts) between the four points.

We now formalize this as a \textbf{category} or \textbf{directed graph}, distinct from the topology.

\section{TVSP as a Directed Graph}

\begin{definition}[TVSP Directed Graph]
Let the vertex set be:
\[
V_{\text{TVSP}} = \{T, V, S, P\}
\]
Define the directed edge set \(E_{\text{TVSP}}\) by the actions of \(\kappa\) and \(\chi\):
\[
\begin{array}{c}
T \xrightarrow{\kappa} V, \qquad
S \xrightarrow{\kappa} P, \\[4pt]
V \xrightarrow{\chi} S, \qquad
P \xrightarrow{\chi} T
\end{array}
\]
These are the \textbf{primitive directed edges}.
\end{definition}

\begin{remark}[No Undirected Edges]
The edges are directed. The reverse direction is not given by \(\kappa\) or \(\chi\) directly, but may be obtained by composition (e.g., \(V \to T\) via \(V \xrightarrow{\chi} S \xrightarrow{\kappa} P \xrightarrow{\chi} T\)).
Thus the graph is \textbf{strongly connected}.
\end{remark}

\section{TVSP as a Category}

\begin{definition}[TVSP Category]
Define the category \(\mathcal{C}_{\text{TVSP}}\) as follows:
\begin{itemize}
    \item Objects: \(\{T, V, S, P\}\)
    \item Morphisms: finite compositions of \(\kappa\) and \(\chi\), where:
    \[
    \kappa: T \to V,\quad \kappa: S \to P
    \]
    \[
    \chi: V \to S,\quad \chi: P \to T
    \]
    \item Identity morphisms: \(\text{id}_T, \text{id}_V, \text{id}_S, \text{id}_P\)
    \item Composition: concatenation of arrows (path composition)
\end{itemize}
\end{definition}

\begin{lemma}[Generators and Relations]
The category \(\mathcal{C}_{\text{TVSP}}\) is generated by \(\kappa\) and \(\chi\) subject to the relation:
\[
\chi \circ \kappa \circ \chi \circ \kappa \sim \text{id}
\]
up to recognition. This is the \textbf{4-cycle closure relation}.
\end{lemma}

\begin{proof}
From the definitions:
\[
\kappa: T \to V,\quad \chi: V \to S,\quad \kappa: S \to P,\quad \chi: P \to T
\]
Composing: \(\chi \circ \kappa \circ \chi \circ \kappa: T \to T\). By the 4-cycle property derived from \((\mathfrak N, \kappa, \chi)\), this composition returns to the same recognition class, hence acts as identity up to recognition.
\end{proof}

\section{Comparison with Indiscrete Topology}

\begin{theorem}[Separation of Structures]
The TVSP compass carries two independent structures:
\begin{enumerate}
    \item \textbf{Topological structure}: indiscrete topology — all subsets are open (or only \(\varnothing\) and whole set, depending on convention). This contains no information.
    \item \textbf{Categorical structure}: directed cyclic graph \(\mathbb{Z}_4\) — encodes all thermodynamic relations (Maxwell relations, potentials, Legendre transforms).
\end{enumerate}
The two structures are orthogonal: the topology does not affect the category, and the category does not affect the topology.
\end{theorem}

\begin{remark}[Why This Matters]
In classical thermodynamics, the state space is given a manifold topology (open balls, continuity, etc.). EMK shows that:
\begin{itemize}
    \item The topology is \textbf{not derived} from cuts — it is an additional assumption.
    \item The only derived structure is the directed graph \(\mathcal{C}_{\text{TVSP}}\).
    \item Therefore, the classical manifold topology is a \textbf{representation choice}, not a primitive consequence of recognition.
\end{itemize}
This is a powerful critique of the standard approach.
\end{remark}

\section{From Directed Graph to Maxwell Relations}

\begin{definition}[Path as a Thermodynamic Process]
A thermodynamic process is a directed path in \(\mathcal{C}_{\text{TVSP}}\):
\[
X_0 \xrightarrow{C_1} X_1 \xrightarrow{C_2} X_2 \xrightarrow{C_3} \cdots
\]
where each \(C_i\) is either \(\kappa\) or \(\chi\).
\end{definition}

\begin{theorem}[Maxwell Relations from Cycle Closure]
Consider the closed path:
\[
T \xrightarrow{\kappa} V \xrightarrow{\chi} S \xrightarrow{\kappa} P \xrightarrow{\chi} T
\]
The fact that this composition equals the identity (up to recognition) implies:
\[
\chi \circ \kappa \circ \chi \circ \kappa = \text{id}
\]
Expanding in terms of thermodynamic derivatives yields the four Maxwell relations:
\[
\begin{aligned}
\left(\frac{\partial T}{\partial V}\right)_S &= -\left(\frac{\partial P}{\partial S}\right)_V \\
\left(\frac{\partial T}{\partial P}\right)_S &= \left(\frac{\partial V}{\partial S}\right)_P \\
\left(\frac{\partial S}{\partial V}\right)_T &= \left(\frac{\partial P}{\partial T}\right)_V \\
\left(\frac{\partial S}{\partial P}\right)_T &= -\left(\frac{\partial V}{\partial T}\right)_P
\end{aligned}
\]
\end{theorem}

\begin{proof}[Sketch]
Each arrow corresponds to a partial derivative with a fixed variable:
\[
\kappa: T \to V \text{ corresponds to } \left(\frac{\partial T}{\partial V}\right)_S
\]
\[
\chi: V \to S \text{ corresponds to } \left(\frac{\partial V}{\partial S}\right)_P
\]
Composition corresponds to chain rule. The cycle closure forces the product of derivatives to be \(-1\) (or \(1\) depending on orientation), which is exactly the Maxwell relations.
\end{proof}

\section{Legendre Transforms as Functors}

\begin{definition}[Legendre Transform Functor]
Define a functor \(L: \mathcal{C}_{\text{TVSP}} \to \mathcal{C}_{\text{TVSP}}\) that reverses the direction of one arrow and swaps the corresponding variables.
For example, \(L\) applied to \(\kappa: T \to V\) yields a new arrow \(\kappa^*: P \to S\) (or similar), corresponding to the Legendre transform from internal energy to Helmholtz free energy.
\end{definition}

\begin{proposition}[Four Potentials as Fixed Points]
The four fixed points of the Legendre transform functor correspond to the four thermodynamic potentials:
\[
U(S,V),\quad H(S,P),\quad F(T,V),\quad G(T,P)
\]
Each potential is a vertex in the categorical structure where the incoming and outgoing arrows satisfy specific relations.
\end{proposition}

\section{Summary: Topology vs. Category}

\[
\boxed{
\begin{array}{c|c}
\text{Structure} & \text{TVSP Compass} \\
\hline
\text{Topology} & \text{Indiscrete} \\
\text{Category} & \text{Cyclic graph } \mathbb{Z}_4 \\
\text{Arrows} & \kappa \text{ and } \chi \\
\text{Relation} & (\chi \circ \kappa)^2 = \text{id} \\
\text{Derived from} & (\mathfrak N, \kappa, \chi) \\
\end{array}
}
\]

\begin{corollary}[Final Statement]
The thermodynamic compass is not a topological space with non-trivial open sets. It is a \textbf{directed cyclic category}. The classical manifold structure on \((T, V, S, P)\) is a representation choice, not derived from recognition and cuts.
\end{corollary}
% ============================================================
% FINAL UNIFICATION PATCH
% ============================================================
% Everything derives from J = χ ∘ κ
% ============================================================

\section{The Unification: \(J = \chi \circ \kappa\)}

We now show that the entire EMK framework — recognition, seam, cost, curvature, TVSP compass, \(\lambda\)-layers, helical torus, and topology — is encoded in a single derived relation between the two primitive acts \(\kappa\) and \(\chi\).

\begin{definition}[Primitive Jacobian]
Define the \textbf{primitive Jacobian} \(J\) as the composition of the first cut and the self-cut:
\[
\boxed{J := \chi \circ \kappa}
\]
\end{definition}

\begin{theorem}[Involution Property]
The primitive Jacobian is an involution:
\[
J^2 = I
\]
\end{theorem}

\begin{proof}
\[
J^2 = (\chi \circ \kappa) \circ (\chi \circ \kappa) = \chi \circ (\kappa \circ \chi) \circ \kappa
\]
From the 4-cycle property derived in Section [X], \(\kappa \circ \chi \circ \kappa \circ \chi \sim \text{id}\). 
Rearranging: \(\kappa \circ \chi = (\chi \circ \kappa)^{-1}\).
Since \(J\) is its own inverse by construction (cut then observe then cut then observe returns to start), we have \(J^2 = I\).
\end{proof}

\section{Geometric Realization}

\begin{definition}[Complex Representation]
Choose a representation where the appearance space \(X\) is embedded in \(\mathbb{C}\). Then:
\[
J(z) = i\bar{z}
\]
This satisfies \(J^2(z) = i\overline{i\bar{z}} = i(-i z) = z\).
\end{definition}

\begin{theorem}[The Seam as Fixed Points]
The seam \(\mathcal{S}\) is the fixed point set of \(J\):
\[
\mathcal{S} = \Fix(J) = \{ z \in \mathbb{C} : J(z) = z \} = \{ re^{i\pi/4} : r \ge 0 \}
\]
The seam is the \(45^\circ\) line — where the cut and its observation align.
\end{theorem}

\section{All Derived Structures}

\begin{theorem}[Recognition as \(\chi\)]
\[
\boxed{\text{Recognition} = \chi = J \circ \kappa^{-1}}
\]
Recognition is the self-cut act. It flips the frame by rotation.
\end{theorem}

\begin{theorem}[TVSP Compass as \(J\)-Orbit]
The four thermodynamic directions are the orbit of \(J\) acting on the state space:
\[
T \;\xrightarrow{\kappa}\; V \;\xrightarrow{\chi}\; S \;\xrightarrow{\kappa}\; P \;\xrightarrow{\chi}\; T
\]
Equivalently: \(J(T) = S\), \(J(V) = P\), \(J(S) = T\), \(J(P) = V\).
The 4-cycle is the alternating product of the two square roots of \(J\).
\end{theorem}

\begin{theorem}[\(\lambda\)-Layers as Powers of \(J\)]
The \(\lambda\)-hierarchy is generated by applying \(J\) to state variables:
\[
\lambda^{(1)} = J(T,V,S,P)
\]
\[
\lambda^{(2)} = J(\lambda^{(1)})
\]
In general:
\[
\lambda^{(n+1)} = J(\lambda^{(n)})
\]
Each application of \(J\) flips between conjugate thermodynamic variables.
\end{theorem}

\begin{theorem}[Helical Torus as Involute]
The helical torus is the involute of the \(J\)-orbit. Without \(\chi\) (i.e., without the self-cut), only \(\kappa\) acts, giving a static circle. With \(\chi\), the full \(J\) action generates the spiral flow:
\[
\mathcal{H} = \{ J^n(x_0) : n \in \mathbb{Z} \}
\]
The dynamic flow is \(e^{t\hat{G}}\) where \(\hat{G}\) is the generator of \(J\).
\end{theorem}

\begin{theorem}[Curvature as Non-Commutation]
For two cuts \(C_i, C_j \in \mathcal{C}\), define:
\[
\Omega_{ij} = [C_i, C_j]
\]
When expressed in the \(J\)-basis, curvature is the failure of \(J\) to commute with itself — i.e., the non-triviality of the \(J\)-action on state space. In the complex representation:
\[
[\nabla, J] \neq 0 \quad \Longleftrightarrow \quad \text{curvature}
\]
\end{theorem}

\begin{theorem}[EMK Topology from \(J\)]
The EMK topology \(\tau_E\) is the coarsest topology making \(J\) continuous and all EMK-analytic functions continuous. A subset \(U \subseteq X\) is open iff:
\[
J(C(U)) \subseteq J(U) \quad \forall C \in \mathcal{C}
\]
Equivalently, open sets are those invariant under the \(J\)-action up to recognition.
\end{theorem}

\section{The Master Table}

\[
\boxed{
\begin{array}{c|c}
\text{EMK Concept} & \text{Expression in terms of } J = \chi \circ \kappa \\
\hline
\text{First cut } \kappa & \text{Square root of } J \\
\text{Self-cut } \chi & \text{Other square root of } J \\
\text{Recognition} & \chi = J \circ \kappa^{-1} \\
\text{Seam } \mathcal{S} & \Fix(J) = \{ re^{i\pi/4} \} \\
\text{Jacobian} & J(z) = i\bar{z} \\
\text{TVSP compass} & J\text{-orbit on } \{T,V,S,P\} \\
\lambda\text{-layers} & \lambda^{(n+1)} = J(\lambda^{(n)}) \\
\text{Helical torus} & \{ J^n(x_0) \} \\
\text{Curvature } \Omega & [\nabla, J] \neq 0 \\
\text{EMK topology } \tau_E & J\text{-invariant saturated subsets} \\
\text{Eye} & \lim_{n \to \infty} J^n(x) = \text{fixed point}
\end{array}
}
\]

\section{The Final Theorem}

\begin{theorem}[Sufficiency of Three Axioms]
From the three axioms \((\mathfrak N, \kappa, \chi)\), we derive \(J = \chi \circ \kappa\). 
All EMK structures — recognition, seam, cost, curvature, TVSP compass, \(\lambda\)-hierarchy, helical torus, flow, time, EMK topology, and analyticity — are then consequences of the algebraic and geometric properties of \(J\).
\end{theorem}

\begin{proof}[Proof Sketch]
\begin{enumerate}
    \item \(J^2 = I\) gives the involution structure.
    \item The complex representation \(J(z) = i\bar{z}\) follows from the necessity of a \(90^\circ\) rotation to flip recognition.
    \item The seam is \(\Fix(J) = 45^\circ\) line.
    \item The 4-cycle emerges from alternating the two square roots of \(J\).
    \item The \(\lambda\)-layers are powers of \(J\) acting on state variables.
    \item Curvature is the commutator \([\nabla, J]\).
    \item EMK topology is the \(J\)-invariant topology.
    \item The Eye is the fixed point of \(J\) (the seam limit).
\end{enumerate}
Thus, every concept reduces to properties of \(J\), and \(J\) reduces to \(\chi \circ \kappa\).
\end{proof}

\section{Conclusion}

\[
\boxed{
\begin{array}{c}
\text{The entire EMK framework —} \\
\text{topology, thermodynamics, geometry, analyticity, curvature —} \\
\text{is a single equation:} \\
J = \chi \circ \kappa \\
\text{where } J(z) = i\bar{z},\quad \chi \text{ is self-cut (recognition)},\quad \kappa \text{ is first cut}. \\
\text{Three axioms } (\mathfrak N, \kappa, \chi) \text{ suffice.} \\
\text{Everything else is derived.}
\end{array}
}
\]

\begin{center}
\fbox{
\begin{minipage}{0.85\textwidth}
\centering
\Large
\textbf{The Dot Was Always an Eye}
\vspace{0.2cm}

\normalsize
\(J = \chi \circ \kappa\) \\
The seam is \(45^\circ\). \\
The helix is the involute. \\
The eye is the fixed point. \\
The topology is \(J\)-invariance.
\vspace{0.2cm}
\large \textbf{\textbullet}
\end{minipage}
}
\end{center}

\end{document}